\uuid{q986}
\titre{Tangentes}
\auteur{hanani ; mzari}
\organisation{exo7}
\licence{CC BY-NC-SA 4.0}
\datecreate{2020-01-20}
\niveau{L1}
\module{Géométrie}
\chapitre{Courbes planes}
\sousChapitre{Courbes paramétrées}
\difficulte{1}
\theme{Courbes paramétrées, tangentes}
\exo7qcmid{lille-2.courbes.001}

\contenu{
\begin{question}
\qclassification{244.01}
\qtags{motcle=tangentes, points doubles}
La trajectoire $\Gamma$ d'une particule en mouvement est donnée par les équations
$$x=t-1\quad \mbox{et} \quad y=t+1.$$
\begin{answers}
\good{$\Gamma$ est la droite d'équation $y=2+x$.}
\bad{La tangente à $\Gamma$ au point $(x(0),y(0))$ est la droite d'équation $y=x$.}
\good{La tangente à $\Gamma$ au point $(x(1),y(1))$ est la droite d'équation $y=2+x$.}
\bad{$\Gamma$ possède un point double.}
\end{answers}
\begin{explanations}
En éliminant le paramètre $t$, on obtient $y=2+x$. $\Gamma$ est la droite d'équation $y=2+x$ et elle est confondue avec sa tangente en n'importe quel point.
\end{explanations}
\end{question}
}
